% FILE: 01_research_question.tex
% PROJECT: ma — M&A Claim Taxonomy and Board Projection Surface

\section{Research Question}
\label{sec:ma-research-question}

\subsection{Problem Statement}
\label{sec:ma-rq-problem}

Mergers and acquisitions are among the highest-stakes capital-allocation decisions in the
enterprise. Yet the evidentiary basis for most M\&A valuations is structurally weak. Synergy
projections are manufactured from spreadsheet extrapolations. Operational health claims derive
from management-prepared KPI decks. Quality-of-Earnings analyses depend on self-reported
financial metrics whose operational underpinnings are unverified. A buyer who pays a merger
premium based on a seller's claim of ``98\% straight-through invoice processing'' has no
independent mechanism to verify that claim before the transaction closes---or to quantify
the discount due if the claim is false.

The problem is not that buyers lack sophistication. It is that the evidentiary
infrastructure for operational claims does not exist in standard M\&A practice. There is no
accepted method by which a seller can make an operational claim that a buyer can
independently replay within a due-diligence timeline. The \texttt{ma/} module asks: can
process mining supply that infrastructure, and if so, what form must the claim taxonomy and
proof protocol take to make M\&A operational claims board-admissible?

\subsection{Old-Regime Assumption Challenged}
\label{sec:ma-rq-old-regime}

The old regime assumes that operational claims in M\&A are inherently soft---that ``process
quality'' is a qualitative judgment captured through management interviews, factory tours, and
sample testing, and that any quantification is an approximation whose margin of error is
accepted as a cost of doing business.

This assumption yields three structural pathologies. First, information asymmetry: the seller
knows the true operational state and the buyer does not, giving the seller an incentive to
present an optimistic picture. Second, latent operational debt: rework intensity, compliance
leakage, SLA penalty exposure, and legacy lock-in are systematically hidden in narrative
presentations and surface only after acquisition, destroying synergy projections. Third,
unverifiable synergy claims: post-merger integration costs are routinely underestimated because
no party has replayed both organizations' processes against a shared model to measure
structural divergence, bottleneck capacity, or behavioral profile similarity.

The v30.1.1 Hostile Takeover Receipt taxonomy (source: \texttt{hostile-takeover-receipts.md})
names this explicitly: ``Without unforgeable HTRs, any PowerPoint synergy claim or denial of
operational debt is legally and financially inadmissible.''

\subsection{Hypothesis}
\label{sec:ma-rq-hypothesis}

\textbf{Hypothesis:} A formal M\&A claim taxonomy grounded in XES/OCEL 2.0 event logs,
alignment-based conformance metrics, and cryptographic receipt chains can make operational
claims board-admissible---enabling independent buyer replication within three days and
compressing traditional seven-week Quality-of-Earnings timelines to automated wasm4pm query
execution.

The hypothesis decomposes into four sub-claims:

\begin{enumerate}
  \item Every material operational claim (EBITDA impact, working capital release, SLA penalty
    exposure, compliance leakage, process standardization) can be expressed as a mathematical
    formula whose inputs are derivable from an OCEL 2.0 or XES event log.
  \item Every such formula can be wrapped in a wasm4pm conformance query, whose execution
    produces a JSON receipt with BLAKE3 log hash, Ed25519 validator signature, and numeric
    fitness/precision results.
  \item A Virtual Data Room (VDR) containing event logs, process models, receipts, and
    provenance lineage files enables an independent buyer's auditor to replay every slide
    assertion in the board presentation.
  \item The Reverse Porter Five Forces inversion is structural: a firm with a multi-year OCEL
    event log stack possesses an evidence moat that is not replicable by a competitor who has
    not preserved equivalent log history, making process evidence a permanent competitive
    barrier.
\end{enumerate}

\subsection{Success Criteria}
\label{sec:ma-rq-success}

The research program defines success as follows:

\begin{itemize}
  \item \textbf{Admissibility thresholds met:} All conformance results satisfy fitness $f \ge
    0.95$, precision $p \ge 0.90$, generalization $g \ge 0.85$, simplicity $s \ge 0.80$, with
    replication tolerance $\Delta f, \Delta p < 10^{-6}$ across independent replay runs.
  \item \textbf{Receipt chain complete:} Every slide in the board projection deck maps to a
    distinct JSON receipt via the Slide-to-Receipt Map, with no orphaned assertions.
  \item \textbf{Auditor path executable:} The five-step auditor verification protocol
    (Reference Extraction $\to$ PROV-O Lineage $\to$ BLAKE3+Ed25519 Integrity $\to$ Soundness
    Proof $\to$ Optimal Alignment Replay) completes without manual intervention.
  \item \textbf{Diligence compression demonstrated:} The seven-step Board Projection Renderer
    pipeline executes end-to-end from claims JSON through sealed PowerPoint artifact within the
    three-day buyer replication window.
\end{itemize}

The checkpoint file \texttt{checkpoint\_\_m\&a-ready\_research\_complete.md} records the
APPROVED verdict for the doctrinal layer. The execution layer---receipt generation from actual
wasm4pm query runs against real enterprise event logs---is deferred (see Chapter~7).
